\documentclass[withoutpreface,bwprint]{cumcmthesis}

% 定义字体 (华文行楷 和 隶书)
\setCJKfamilyfont{hwxk}[AutoFakeBold]{STXingkai} % 华文行楷
\newcommand{\huawenxingkai}{\CJKfamily{hwxk}}
\setCJKfamilyfont{mylisu}[AutoFakeBold]{LiSu} % 隶书
\newcommand{\mylishu}{\CJKfamily{mylisu}}

\begin{document}

\begin{titlepage}
    \centering
    \vspace*{0.5cm}
    \includegraphics[width=12cm]{figures/图标.png} \par
    \vspace{1.0cm}
    % {\huawenxingkai \fontsize{36pt}{\baselineskip}\selectfont \textbf{智能制造技术} \par}
    \vspace{1.2cm}
    {\mylishu \fontsize{60pt}{\baselineskip}\selectfont 商业伦理与\\会计职业\par}
    \vspace{1.5cm}
    {\heiti \fontsize{36pt}{\baselineskip}\selectfont \textbf{结\hspace{0.5em}课\hspace{0.5em}作\hspace{0.5em}业} \par}
    \vspace{3.0cm}
    \begin{center}
        \zihao{3}
        \renewcommand\arraystretch{1.5}
        \begin{tabular}{c c}
            \makebox[4.5em][s]{题\hfill 目} & \underline{\makebox[21em][c]{}} \\
            %  & \underline{\makebox[21em][c]{}} \\
            \makebox[4.5em][s]{专\hfill 业} & \underline{\makebox[21em][c]{}} \\
            \makebox[4.5em][s]{班\hfill 级} & \underline{\makebox[21em][c]{}} \\
            \makebox[4.5em][s]{学\hfill 生} & \underline{\makebox[21em][c]{}} \\
            % \makebox[4.5em][s]{指导教师}   & \underline{\makebox[21em][c]{郭正~杨新刚}} \\
        \end{tabular}
    \end{center}
    \vfill
    % {\songti \zihao{4} 2026年1月10日}
    \vspace{1cm}
\end{titlepage}

\newpage
\thispagestyle{empty}
\vspace*{-2cm} % 调整图片上边距
\begin{center}
    \includegraphics[width=8cm]{figures/图片1.png}
\end{center}
\vspace{0.5cm} % 调整图片与标题的间距
\begin{center}
    {\heiti \zihao{2} 智能制造技术课程设计任务书}
\end{center}

\vspace{1.5cm}

\noindent {\heiti \zihao{4} 设计题目：} \underline{\makebox[10em][c]{尾座体}} \zihao{4} 零件机械加工工艺规程 \par
\vspace{0.8em}
\noindent \hspace{5.5em} 及 \underline{\makebox[14em][c]{尾座体加工斜孔}} 工序专用夹具设计

\vspace{1.0cm}

\noindent {\heiti \zihao{4} 设计要求：}

\vspace{0.5em}
{\zihao{4} \renewcommand\arraystretch{1.5}
\noindent 1、制订零件机械加工工艺，填写机械加工工艺过程卡一张，设计夹具工序的工序卡一张；

\noindent 2、设计指定工序的专用夹具一套，绘制夹具装配图一张；

\noindent 3、编写机械制造工艺课程设计说明书一份（4000$\sim$6000字）。
}

\vspace{1.0cm}

\noindent {\heiti \zihao{4} 设计内容：}

\vspace{0.5em}
{\zihao{4} \renewcommand\arraystretch{1.5}
\noindent 1、熟悉及审查零件，选择毛坯种类及制造方法，绘制零件毛坯联合图1张或零件图与毛坯图各1张；

\noindent 2、按设计任务书的要求，对零件进行工艺分析，拟定工艺方案；

\noindent 3、拟定零件的机械加工工艺过程，选择各道工序的加工设备和工艺装备（刀具、夹具、量具、辅具），确定重要工序切削用量及工序尺寸，计算某一代表工序的工时定额。

\noindent 4、填写工艺文件：工艺过程卡一套、设计夹具工序的工序卡一份；

\noindent 5、设计重要工序中的一种专用夹具（车、铣、钻、镗等夹具），绘制装配图一张；

\noindent 6、撰写设计说明书一份。
}

\vfill

\begin{flushright}
    \zihao{4}
    \renewcommand\arraystretch{1.5}
    \begin{tabular}{l c}
        姓名：& \underline{\makebox[8em][c]{安家乐}} \\
        学号：& \underline{\makebox[8em][c]{3230212006}} \\
        班级：& \underline{\makebox[8em][c]{工程232}} \\
    \end{tabular}
\end{flushright}
\vspace{2cm}

\newpage


\pagenumbering{Roman}
\tableofcontents
\newpage
\pagenumbering{arabic}
\setcounter{page}{1}

% 增加段落间距，改善阅读体验并增加篇幅
\setlength{\parskip}{0.8ex plus 0.2ex minus 0.1ex}

\section{序言}

\subsection{设计的背景与意义}
智能制造课程设计作为我们机械类专业学生极为关键的综合性实践教学环节，起着承前启后的重要作用。它不仅是对过去所学理论知识的一次全面检验，更是连接理论学习与工程实践的桥梁，对于我们将书本知识转化为实际工程能力具有深远的意义。

在当今高度发达的现代制造业体系中，工艺规程占据着核心地位，它是指导现场生产、组织生产流程以及科学管理生产活动的主要技术纲领性文件。一份科学、严谨且合理的工艺规程，其价值不仅体现在能够稳定地保证产品加工质量，规避生产风险，更在于它能有效优化生产节奏，显著提高生产效率，并从源头上降低制造成本，提升企业的市场竞争力。与此同时，专用夹具的设计与应用，则是落实工艺规程、确保零件加工精度稳定性不可或缺的物质基础和技术手段。

本次设计的核心对象——机床尾座体，虽然从外观上看结构相对紧凑，但实则涵盖了箱体类零件加工过程中的诸多典型特征与技术难点，例如：复杂平面的铣削加工、精密孔系的镗铰加工、以及对位置精度（如平行度、垂直度、同轴度）的严格控制等。通过对该典型零件的完整工艺设计流程，我通过查阅资料、计算校核等方式，旨在达到以下目的：
\begin{enumerate}
    \item \textbf{理论深化与应用}：将课堂上学到的诸如“六点定位原理”、“基准重合原则”、“误差复映规律”等抽象的理论概念，具体应用到真实的工程实例分析中，实现理论与实践的深度融合。
    \item \textbf{工程能力培养}：通过查找和使用《机械加工工艺手册》等专业工具书，培养独立检索资料、分析复杂工程问题以及解决实际生产问题的综合能力。
    \item \textbf{工程素养养成}：在设计过程中，逐步建立起“质量为本、效率优先、成本控制”三位一体的现代工程经济观念，并养成严谨、规范、细致的工程设计习惯，为未来的职业生涯打下坚实基础。
\end{enumerate}

\subsection{设计任务分解}
为了确保本次课程设计能够有序、高质量地完成，主要任务被科学地分解为以下三个递进的阶段：
\begin{itemize}
    \item \textbf{阶段一：工艺审查与毛坯确定。} 深入分析图纸技术要求，全面判断零件结构的工艺性，合理选择毛坯的制造种类及确定加工余量。
    \item \textbf{阶段二：工艺规程设计。} 这是本次设计的核心部分，涵盖了定位基准的科学选择、工艺路线的逻辑拟定、详细的工序尺寸及公差计算、以及切削用量及工时定额的精确确定。
    \item \textbf{阶段三：专用夹具设计。} 针对加工过程中的关键瓶颈工序（如斜面孔加工），设计一套切实可行的专用夹具，包括定位方案的构思、夹紧机构的设计、精细的误差分析以及装配图的绘制。
\end{itemize}


\section{零件分析}

\subsection{零件的作用分析}
机床尾座体作为机床总体结构中的重要组成部分，广泛应用于各类普通车床（如经典的 CA6140 型车床）以及各类高精度外圆磨床中。它通常被精确地安装在床身的导轨之上，物理位置上位于主轴箱的直接对面，与主轴箱协同工作。通过对零件图的深入剖析，我们可以将其具体功能详尽地归纳为以下几个核心方面：

\subsubsection{关键的支撑与刚性增强功能}
尾座体的核心功能区在于其上部的 $\Phi17$H6 主轴孔，该孔内部精密配合安装有可以轴向移动的顶尖套筒（通常称为尾座套筒）。在实际的车削或磨削加工过程中，顶尖前端紧密地插在工件端面的中心孔内。
\begin{itemize}
    \item \textbf{辅助支承}：对于一般的轴类零件，它常与主轴卡盘配合，形成标准的“一夹一顶”装夹方式。
    \item \textbf{刚性补强}：特别是针对长径比较大的细长轴类工件，如果不使用尾座顶尖进行辅助支撑，工件极易在巨大的径向切削力作用下产生弯曲变形，导致出现“让刀”现象，严重影响加工精度和表面质量。尾座顶尖的介入，显著增加了工件工艺系统的整体刚性，从而确保了加工过程的稳定性。
\end{itemize}

\subsubsection{进给与孔加工拓展功能}
除了这一基础的支撑作用外，尾座套筒内部设计的莫氏锥孔结构，提供了极大的工具兼容性。它可以直接安装麻花钻、铰刀、中心钻等各类孔加工标准刀具。
操作者通过手动旋转尾座手轮，经由丝杠螺母机构带动套筒作轴向直线运动，从而实现刀具的轴向进给。这极大地拓展了车床的加工能力，使其不仅能胜任外圆车削，还能高效、方便地完成工件端面的中心孔钻削、深孔钻削作业以及扩孔、铰孔等复杂的孔系加工任务。

\subsubsection{高精度的导向与定位功能}
尾座体底部的导轨槽（设计图中通常呈现为燕尾槽或矩形槽结构）是其与机床床身导轨进行精密配合的关键界面。这一配合面的制造精度直接决定了尾座主轴中心线与床头箱主轴中心线在机床全长范围内的重合度（即同轴度指标）。
如果尾座体在移动过程中发生微量的偏斜或中心高度出现偏差，加工出的圆柱体工件将不可避免地产生锥度误差（如出现“马鞍形”或“鼓形”缺陷），这将直接判定零件为废品。因此，底面和导轨面的加工质量是保证机床几何精度的生命线。

\subsection{零件的工艺性审查}
依据零件图样，我对尾座体的结构特点和技术指标进行了详细的工艺性审查。

\subsubsection{结构特点分析}
该零件属于典型的不规则类箱体零件，其结构工艺性具有以下特点：
\begin{enumerate}
    \item \textbf{壁厚不均，易产生变形}：零件既有厚实的底座部分，又有相对较薄的侧壁及凸台。这种壁厚差异在铸造冷却过程中极易导致散热速度不同，从而产生较大的残余内应力。如果处理不当，后续加工中会发生翘曲变形。因此，工艺路线安排中必须包含充分的“人工时效处理”工序。
    \item \textbf{多面加工，基准转换复杂}：零件需要加工底面、顶面、前端面、后端面、侧面等多个平面，且这些平面之间存在严格的垂直度或平行度要求。
    \item \textbf{孔系空间位置复杂}：包含 $\Phi17$ 主轴孔和 $\Phi14$ 侧孔，且两孔轴线在空间上呈现交叉关系，非正交加工增加了装夹和对刀的难度。
\end{enumerate}

\subsubsection{关键技术要求解读}
我对图纸上标注的几项核心公差及其加工难点进行了深入分析：

\textbf{1. $\Phi17$H6 主轴孔：}
\begin{itemize}
    \item \textbf{尺寸精度（H6级）}：这是极高的尺寸精度要求。通常情况下，H7 级已属于精密孔加工范畴，而 H6 级的公差带极窄（约为 $0.011$mm）。这意味着加工过程中刀具磨损、机床主轴跳动等微小误差都可能导致超差。
    \item \textbf{表面粗糙度（Ra 0.4 $\mu m$）}：这已接近研磨或精密磨削的表面质量标准。普通的钻削工艺 $Ra$ 值通常在 12.5 以上，标准铰削一般也只能达到 $Ra 1.6$。要稳定达到 $Ra 0.4$，必须采用多道次的精细铰削工艺（如浮动精铰），并配合优良的切削液进行润滑和冷却。
\end{itemize}

\textbf{2. 底面导轨面：}
\begin{itemize}
    \item \textbf{表面粗糙度（Ra 1.6 $\mu m$）}：底面是与床身导轨接触的滑动摩擦面，较好的表面光洁度能有效减少磨损，延长机床使用寿命，同时也影响着移动时的平稳性。
\end{itemize}



\section{工艺规程制订}

\subsection{毛坯的确定}
\subsubsection{材料选择的科学依据}
经过对设计图纸及工况的详细查阅，确定零件材料选用 HT200（灰铸铁）。这一选择是基于以下几方面的综合考量：
\begin{itemize}
    \item \textbf{优良的减震性能}：机床尾座在切削工作时主要承受复杂的压力载荷和切削引起的频率振动。灰铸铁金相组织中存在的片状石墨能有效割裂金属基体，从而切断振动波的传播路径，具有钢材无法比拟的极好消振性（减震性），这有利于保证机床的加工精度。
    \item \textbf{出色的耐磨与润滑性}：石墨本身具有润滑作用，脱落后留下的凹坑可以储存润滑油。HT200 作为导轨接触面材料，能有效降低摩擦系数，通过人工刮研后能获得极佳的导轨接触性能，延长使用寿命。
    \item \textbf{良好的切削与铸造工艺性}：HT200 熔点较低，流动性好，收缩率小，适合铸造形状复杂的薄壁壳体零件。同时，其切削时产生的是崩碎切屑，断屑容易，不缠绕刀具。
    \item \textbf{经济性原则}：相较于铸钢或球墨铸铁，灰铸铁成本低廉，货源充足，非常适合作为中低载荷机床床身、箱体及尾座的制造材料。
\end{itemize}

\subsubsection{毛坯制造形式的选择}
\begin{itemize}
    \item \textbf{工艺甄选}：鉴于零件拥有不对称的复杂外观，且包含较多的凸台、内腔结构，若采用锻造工艺难以成型复杂内腔；若直接从型材进行全切削加工，则会导致材料利用率极低，且加工工时成倍增加。综上所述，选用\textbf{铸造}工艺是唯一合理的选择。
    \item \textbf{铸造方法}：结合生产纲领分析，该零件的生产类型属于中批量生产模式。零件轮廓尺寸中等，对毛坯表面的光洁度及尺寸精度要求一般，故选用技术成熟、成本可控的\textbf{砂型铸造}工艺。
    \item \textbf{毛坯结构设计}：
    \begin{itemize}
        \item 针对 $\Phi17$ 主轴孔：由于该孔直径相对较小，若强行在毛坯上铸出该孔，所需的型芯不仅细长，而且容易在浇注冲击下产生位移（被冲歪）或发生烧结现象。此外，铸孔内表面形成的硬皮（冷硬层）会极大加剧后续钻头切削刃的磨损，甚至导致钻头引偏，影响孔的直线度。综合考虑，决定\textbf{不铸出 $\Phi17$ 孔}，采用实心铸造方案。
        \item 针对 $\Phi35$ 圆柱外形：由于尺寸较大，利于造型，故可以铸出。
        \item 余量设定：参考同类铸件标准，预留出 2-3mm 的单边机械加工余量。
    \end{itemize}
\end{itemize}

\subsection{定位基准的选择与论证}
基准选择是制定工艺规程中最关键、最核心的环节之一。基准选择的正确与否，直接关系到零件加工质量的优劣、生产效率的高低以及夹具结构的复杂程度。如果基准选择不当，即使使用再高精度的机床设备，也难以加工出合格的产品，甚至造成大量废品。

\subsubsection{粗基准的选择原则与方案}
粗基准主要用于零件加工的起始工序。它的核心作用有两个：一是建立加工表面与不加工表面之间的正确位置关系；二是为后续工序提供可靠的精基准。
\begin{itemize}
    \item 备选方案 A：以顶面凸台为粗基准。
    \item 备选方案 B：以 $\Phi35$ 圆柱体的外圆表面为粗基准。
    \item \textbf{方案论证与决策}：
    依据《机械加工工艺手册》中的粗基准选择原则：“若零件上有不加工表面，应优先以不加工表面为粗基准”。虽然本零件的顶面和圆柱外圆后续均需进行切削加工，但我们的核心目标是保证将来加工出来的 $\Phi17$ 主轴孔中心线能位于 $\Phi35$ 圆柱体的正几何中心（即保证加工后的壁厚均匀一致）。
    若偏离中心太大，可能导致壁厚过薄而强度不足。因此，应选择与孔同轴的圆柱外圆作为粗基准。
    \item \textbf{结论}：最终决定选择 $\Phi35$ 圆柱体外圆（配合侧面）作为第一道工序的粗基准，主要用于加工底面。
\end{itemize}

\subsubsection{精基准的选择原则与方案}
精基准的选择直接决定了零件最终的加工精度。遵循指导书第二章原则，我做了如下考量：
\begin{enumerate}
    \item \textbf{基准重合原则}：
    零件图样的设计基准是底面（导轨面）和 $\Phi17$ 孔中心线。在加工 $\Phi17$ 孔及其它主要表面时，若直接以底面作为定位基准，则定位基准与设计基准完全重合，理论上的基准不重合误差 $\Delta_B = 0$。这是保证加工精度的最理想状态。
    \item \textbf{基准统一原则}：
    观察零件图可知，零件上绝大多数表面（如前后端面、顶面、侧孔等）的设计基准都是底面。如果我们在加工这些表面时，都统一长期使用底面作为定位基准，就能有效避免因多次基准转换带来的尺寸链累积误差。同时，这种做法可以简化工艺装备的设计，使得多道工序可以共用同一套定位方案或通用夹具。
\end{enumerate}
\textbf{结论}：确定选用加工后的\textbf{底面}和\textbf{导轨槽侧面}作为全过程的主要精基准。

\subsection{工艺路线的制定}

\subsubsection{关键表面加工方案的分析}
根据《机械制造课程设计指导书》第三章的经济加工精度表，我对关键表面进行了深入的方案比选。

1. 底面加工方案（目标：$Ra 1.6$）
\begin{itemize}
    \item 方案一：粗铣 - 半精铣。
    查表 3-9 可知，半精铣的经济精度通常为 IT8 $\sim$ IT10，粗糙度 $Ra 2.5 \sim 6.3 \mu m$。此方案难以稳定达到 $Ra 1.6$ 的要求，可能会导致基准不平。
    \item 方案二：粗铣 - 精铣 - 磨削。
    磨削可以轻松达到 $Ra 0.8$ 甚至更低。但磨削效率低，且对于铸铁件底面，若无特殊硬度要求，一般不推荐磨削。
    \item 方案三：粗铣 - 精铣。
    查表可知，精铣的经济粗糙度范围是 $Ra 0.63 \sim 5 \mu m$。只要选用合理的切削用量（高转速、小进给）和硬质合金刀具，完全可以达到 $Ra 1.6$。此方案效率高、成本低。
    \item 结论：选用 方案三（粗铣-精铣）。
\end{itemize}

2. $\Phi17$H6 主轴孔加工方案（目标：$Ra 0.4$, H6）
\begin{itemize}
    \item 方案一：钻 - 扩 - 铰。
    这是最常用的孔加工路线。查表 3-2，该路线经济精度为 IT8 $\sim$ IT9，粗糙度 $Ra 1.6$。
    \textit{思考：} 这显然无法满足 H6 和 $Ra 0.4$ 的高要求。如果按此方案，必须完全依赖工人的操作技巧，废品率极高。
    \item 方案二：钻 - 扩 - 粗铰 - 精铰。
    查表 3-8，精铰（或手铰）可以达到 IT6 $\sim$ IT7 级精度，粗糙度可达 $Ra 0.16 \sim 0.8 \mu m$。
    \textit{思考：} 增加一道精铰工序，可以专门用于纠正粗铰后的微量误差，并修光孔壁。这是实现高精度的必要手段。
    \item 方案三：镗孔。
    镗孔主要用于大孔或旋转体。对于 $\Phi17$ 的小孔，镗刀杆刚性差，易震颤，难以保证质量。
    \item 结论：选用 方案二（钻-扩-粗铰-精铰）。
\end{itemize}

3. 侧面斜孔加工方案（目标：$Ra 12.5$）
\begin{itemize}
    \item 方案一：钻孔。
    查表可知，普通钻孔的经济粗糙度范围是 $Ra 12.5 \sim 25 \mu m$。由于设计要求仅为 $Ra 12.5$，且该孔主要用于穿过油路或连接，非精密配合孔，因此一道钻孔工序即可满足要求。
    \item 方案二：钻 - 扩。
    扩孔可降低粗糙度至 $Ra 6.3$ 左右，但工序增加，成本上升，且对于 $Ra 12.5$ 的要求属于过度加工。
    \item 结论：选用 方案一（一次钻孔成型）。并且需要设计专用的角度夹具来保证斜孔的角度。
\end{itemize}




\subsubsection{工序顺序的安排逻辑}
根据指导书第二章第五节“加工顺序安排的原则”，我进行了如下编排：
\begin{enumerate}
    \item 基准先行：
    既然底面是精基准，那么必须第一步就把底面加工好。所以工序 20、25 安排铣底面。
    \item 先面后孔：
    \textit{思考：} 为什么要先面后孔？如果先钻 $\Phi17$ 孔，再铣端面，那么钻头切入和切出时面对的是粗糙不平的铸造表面，极易引偏，打断钻头，且孔口会有毛刺。
    \textit{决策：} 必须先将前后端面铣平（工序 45），再进行钻孔（工序 65）。这样钻头是在平整的表面切入，受力均匀，孔位准确。
    \item 先粗后精：
    为了消除切削热和内应力对精度的影响，粗加工（如粗铣、钻孔）应在前，精加工（精铣、精铰）应在后。
    特别注意：在粗加工（铣底面、钻孔）之后，精加工之前，我特意安排了一道“中间检验”工序，防止将废品流入精加工工序，浪费工时。
\end{enumerate}

\subsubsection{最终工艺路线确定}
经过上述反复推敲，确定的工艺路线如下：
\begin{itemize}
    \item 工序05 铸造：砂型铸造。
    \item 工序10 时效：人工时效，消除内应力。
    \item 工序15 涂漆。
    \item 工序20 粗铣底面：以外圆定位，留余量。
    \item 工序25 精铣底面：保证 $Ra 1.6$，形成统一精基准。
    \item 工序30 中间检验：关键点，确认基准是否合格。
    \item 工序40 铣导轨槽及侧面。
    \item 工序45 铣前后端面：为钻孔做准备。
    \item 工序50 铣顶面凸台。
    \item 工序55 精铣左端斜面：需专用角度夹具。
    \item 工序60 铣 $\Phi35$ 圆柱外圆。
    \item 工序65 钻 $\Phi17$ 孔至 $\Phi15$：粗加工。
    \item 工序70 扩 $\Phi17$ 孔至 $\Phi16.8$：半精加工，纠正轴线。
    \item 工序75 粗铰 $\Phi17$ 孔至 $\Phi16.94$：精加工前奏。
    \item 工序80 精铰 $\Phi17$H6 孔：最终保证 $Ra 0.4$ 和 H6。
    \item 工序85 钻 $\Phi14$ 孔，锪Φ26×1的沉头孔：侧孔精度要求低，一次钻成。
    \item 工序90 钳工：攻丝、去毛刺。
    \item 工序105 终检。
\end{itemize}

最终汇总如下表所示（表 \ref{tab:process_scheme} 和表 \ref{tab:process_stage}）。

\begin{table}[H]
    \centering
    \caption{尾座体主要表面加工方案汇总表}
    \label{tab:process_scheme}
    \renewcommand{\arraystretch}{1.3} % 增加行高
    \begin{tabular}{l c c l l}
        \toprule
        \textbf{加工表面名称} & \textbf{精度等级} & \textbf{Ra ($\mu m$)} & \textbf{加工方案 (路线)} & \textbf{备注 (依据)} \\
        \midrule
        底面 (导轨配合面) & IT8 & 1.6 & 粗铣 $\rightarrow$ 精铣 & 依据表3-9，精铣满足要求 \\
        $\Phi17$H6 主轴孔 & IT6 & 0.4 & 钻 $\rightarrow$ 扩 $\rightarrow$ 粗铰 $\rightarrow$ 精铰 & 依据表3-8，需精铰达 IT6 \\
        $\Phi14$ 侧孔 & IT11 & 12.5 & 一次钻孔 & 钻削即可满足 IT11 \\
        左端斜面 ($45^\circ$) & IT8 & 1.6 & 粗铣 $\rightarrow$ 精铣 & 需专用角度夹具 \\
        导轨槽侧面 & IT9 & 6.3 & 粗铣 $\rightarrow$ 半精铣 & 保证槽宽及对称度 \\
        $\Phi35$ 圆柱外表面 & IT9 & 6.3 & 粗铣 $\rightarrow$ 半精铣 & 粗基准表面 \\
        前后端面 & IT11 & 6.3 & 粗铣 & 保证垂直度，为钻孔做准备 \\
        $\Phi26$ 沉头孔 & IT12 & 12.5 & 锪孔 & / \\
        M6 螺纹孔 & 7H & / & 钻底孔 $\rightarrow$ 攻丝 & / \\
        \bottomrule
    \end{tabular}
\end{table}


\begin{table}[H]
    \centering
    \caption{工艺过程加工阶段划分表}
    \label{tab:process_stage}
    \renewcommand{\arraystretch}{1.3}
    \begin{tabular}{l p{5cm} p{6cm}}
        \toprule
        \textbf{加工阶段} & \textbf{主要工序内容} & \textbf{阶段目的与说明} \\
        \midrule
        \textbf{1. 粗加工阶段} & 
        \begin{itemize}
            \item 粗铣底面 (工序20)
            \item 粗铣导轨槽及侧面
            \item 粗铣前后端面、顶面
            \item 铣 $\Phi35$ 圆柱外圆
        \end{itemize} & 
        1. 切除毛坯大部分余量，提高生产率；\newline
        2. 及时发现铸造气孔、砂眼等缺陷；\newline
        3. 为半精加工提供定位基准。 \\
        \hline
        
        \textbf{2. 半精加工阶段} & 
        \begin{itemize}
            \item \textbf{精铣底面} (工序25)
            \item 钻、扩 $\Phi17$ 孔
            \item 钻 $\Phi14$ 孔及锪孔
        \end{itemize} & 
        1. \textbf{确立精基准}：精铣底面作为后续关键工序的统一基准；\newline
        2. 修正粗加工后的位置误差，为精加工做准备；\newline
        3. 完成次要表面的加工（如侧孔）。 \\
        \hline
        
        \textbf{3. 精加工阶段} & 
        \begin{itemize}
            \item 粗铰、\textbf{精铰 $\Phi17$H6 孔}
            \item 精铣左端 $45^\circ$ 斜面
        \end{itemize} & 
        1. 保证关键表面（主轴孔）达到最终的尺寸精度 (IT6) 和表面粗糙度 (Ra 0.4)；\newline
        2. 确保孔轴线对底面的平行度要求。 \\
        \bottomrule
    \end{tabular}
\end{table}

\subsection{机械加工余量、工序尺寸及公差的确定}

据计算和查表可知，3mm机械加工余量已足够，取零件各表面的机械加工总余量均为3mm，以加工 $\Phi14$ 斜孔的工序为例，说明工序尺寸及公差的确定过程。

本工序内容：钻 $\Phi14$ 孔，锪 $\Phi26 \times 1$ 沉头孔。
加工设备：X51 立式铣床。
特征分析：根据参考资料，该孔轴线与底面成 $45^\circ$ 夹角（或位于斜面上），需采用角度夹具加工。

\begin{enumerate}
    \item 工步 1：钻 $\Phi14$ 孔
    \begin{itemize}
        \item 余量：实心钻出，双边余量 $2Z = 14$ mm。
        \item 工序尺寸：$\Phi14$ mm。
        \item 公差：按 IT12 级控制，$\Phi14^{+0.18}_{0}$ mm。
    \end{itemize}
    \item 工步 2：锪 $\Phi26$ 沉头孔
    \begin{itemize}
        \item 工序尺寸：$\Phi26$ mm，深度 1 mm。
    \end{itemize}
\end{enumerate}

\subsection{切削用量及工时定额的计算}

\subsubsection{工步 1：钻 $\Phi14$ 孔}
\begin{enumerate}
    \item 背吃刀量：$a_p = d/2 = 7$ mm。
    \item 进给量 $f$：
    参考指导书及网络资料验证：对于 HT200 材料，$\Phi14$ 钻头在刚性较好的铣床上加工，进给量可取 $0.35 \sim 0.45$ mm/r，
    取 $f = 0.4$ mm/r。
    \item 切削速度 $v_c$：
    查资料及指导手册，取 $v_c \approx 21$ m/min。
    \item 主轴转速 $n$：
    \[ n = \frac{1000 v_c}{\pi d} = \frac{1000 \times 21}{3.14 \times 14} \approx 477 \text{ r/min} \]
    查 X51 说明书，选取 $n = 475$ r/min（与工序卡上的 470 接近）。
    \item 机动时间 $T_{m1}$：
    孔深 $L=30$mm（斜向深度），切入切出 $L_1+L_2=5$mm。
    \[ T_{m1} = \frac{30 + 5}{475 \times 0.4} \approx 0.18 \text{ min} \]
\end{enumerate}

\subsubsection{工步 2：锪 $\Phi26$ 沉头孔}
\begin{enumerate}
    \item 进给量：手动或低速进给，取 $f = 0.2$ mm/r。
    \item 转速：取 $n = 150$ r/min。
    \item 切削速度 $v_c$：
    \[ v_c = \frac{\pi d n}{1000} = \frac{3.14 \times 26 \times 150}{1000} \approx 12.25 \text{ m/min} \]
    \item 切削深度 $a_p$：
    \[ a_p = \frac{D-d}{2} = \frac{26-14}{2} = 6 \text{ mm} \]
    \item 机动时间：$T_{m2} = \frac{1+2}{150 \times 0.2} = 0.1 \text{ min}$。
\end{enumerate}

\begin{table}[h]
    \centering
    \caption{工序切削用量计算结果汇总}
    \renewcommand\arraystretch{1.5}
    \begin{tabular}{|c|c|c|c|c|}
        \hline
        \textbf{工步内容} & \textbf{主轴转速} & \textbf{切削速度} & \textbf{进给量} & \textbf{切削深度} \\
         & (r/min) & (m/min) & (mm/r) & (mm) \\
        \hline
        工步1：钻 $\Phi14$ 孔 & 475 & 20.9 & 0.4 & 7 \\
        \hline
        工步2：锪 $\Phi26$ 孔 & 150 & 12.25 & 0.2 & 6 \\
        \hline
    \end{tabular}
\end{table}

单件工时：$T = T_{m1} + T_{m2} + T_{aux} \approx 0.18 + 0.1 + 0.8 \approx 1.08 \text{ min}$。

\section{专用夹具设计（45度斜度钻模）}

\subsection{设计问题的提出与难点分析}
本工序的加工任务是在 X51 立式铣床上钻削加工 $\Phi14$ 侧孔。
\begin{itemize}
    \item \textbf{特殊性分析}：根据零件结构图纸的几何关系分析，$\Phi14$ 孔的轴线与基准底面并不垂直，而是呈现 $45^\circ$ 的夹角。常规的立式铣床主轴只能做垂直升降运动，无法直接加工斜孔。
    \item \textbf{解决方案}：为了解决这一难题，必须设计一套带有 \textbf{45度斜度底座} 的专用钻孔夹具。其核心设计思想是“空间旋转”，即通过将工件倾斜 $45^\circ$ 安装在夹具上，使得被加工孔原本倾斜的轴线转变为与地面垂直，从而利用立式铣床主轴的垂直进给功能完成加工。
\end{itemize}

\subsection{定位方案的详细设计}
根据机械制造工艺学中的“六点定位原理”，为了保证加工精度，工件在夹具中必须具有确定的位置，即需合理限制其自由度。本夹具的定位方案设计如下：

\begin{enumerate}
    \item \textbf{主定位基准（底面 A）的设计}：
    选用一块精密加工的\textbf{45度大平面定位板}，平稳安装在夹具体的主体上。
    \begin{itemize}
        \item \textit{限制自由度分析}：限制了 3 个自由度，分别是沿 $Z'$ 轴的移动、绕 $X'$ 轴的旋转、以及绕 $Y'$ 轴的旋转。
        \item \textit{功能作用}：这是最重要的定位元件，直接保证 $\Phi14$ 孔轴线与底面之间的严格角度关系。
    \end{itemize}

    \item \textbf{导向定位基准（侧面 B）的设计}：
    选用标准的\textbf{侧面固定座}，紧贴在斜面的侧边，作为一个长导向定位元件。
    \begin{itemize}
        \item \textit{限制自由度分析}：限制了 2 个自由度，分别是沿 $X'$ 轴的移动、绕 $Z'$ 轴的旋转。
        \item \textit{功能作用}：保证孔相对于侧面的位置度，防止加工时工件发生侧向偏移。
    \end{itemize}

    \item \textbf{止推定位基准（前端面）的设计}：
    选用一枚高强度的\textbf{固定挡销}安装在斜面底部。
    \begin{itemize}
        \item \textit{限制自由度分析}：限制了 1 个自由度，即沿 $Y'$ 轴（下滑方向）的移动。
        \item \textit{功能作用}：既起到定位作用，确定轴向位置；又起到承受切削力和重力沿斜面下滑分力的止推作用。
    \end{itemize}
\end{enumerate}

\subsection{对刀与夹紧设计}
\begin{itemize}
    \item 对刀装置：在夹具斜面上方设置钻套。加工前，工人用塞尺调整铣床主轴（Z轴及XY轴）位置，确保钻头中心与斜放后的孔中心重合。
    \item 夹紧机构：采用 移动式螺旋压板。
    \item 夹紧方向：垂直于 $45^\circ$ 斜面，压向大平面定位板。
\end{itemize}

\subsection{切削力与夹紧力的详细计算}

\subsubsection{1. 切削力理论计算}
为了设计出安全可靠的夹具，必须准确计算切削过程中产生的最大轴向力。查阅《机械制造课程设计指导书》表 3-28 关于钻削加工的相关数据，对于硬度为 190HBS 的灰铸铁材料，其钻孔轴向力 $F_f$ 的计算经验公式如下：
\[ F_f = C_F d_0^{z_F} f^{y_F} k_F \]
式中各参数的意义及取值如下：
\begin{itemize}
    \item $C_F$：系数，取决于工件材料和刀具类型。查表得 $C_F = 610$。
    \item $d_0$：钻头直径。本工序钻头直径 $d_0 = 14$ mm。
    \item $z_F$：直径指数。查表得 $z_F = 1.0$。
    \item $f$：进给量。本工序选用 $f = 0.4$ mm/r。
    \item $y_F$：进给量指数。查表得 $y_F = 0.8$。
    \item $k_F$：加工材料硬度修正系数。$k_F = (HB/190)^{0.6} = (190/190)^{0.6} = 1.0$。
\end{itemize}

将上述数据代入公式进行计算：
\[ 
\begin{aligned}
F_f &= 610 \times 14^{1.0} \times 0.4^{0.8} \times 1.0 \\
    &= 610 \times 14 \times 0.480 \\
    &\approx 4099 \text{ N} 
\end{aligned}
\]
由此可知，钻孔时的最大垂直轴向力约为 4099 N。

\subsubsection{2. 所需夹紧力的估算}
在本工序中，工件倾斜 45 度放置，切削力 $F_f$ 的方向是垂直向下压紧工件的，这在一定程度上有利于定位的稳定性。然而，钻削过程中的扭矩以及钻头切入切出时的振动，仍可能导致工件松动。此外，工件重力 $G$ 会产生一个沿斜面下滑的分力 $G \sin 45^\circ$。

为了确保绝对安全，我们设定较大的安全系数 $K$。参考机床夹具设计手册，取安全系数 $K = 2.5$。则理论所需夹紧力 $W_k$ 为：
\[ W_k = K \times F_f = 2.5 \times 4099 \approx 10247 \text{ N} \]
夹具必须时刻提供大于 10247 N 的夹紧力才能保证加工过程万无一失。

\subsubsection{3. 夹紧机构的强度校验}
根据上述计算结果，我们需要一个能稳定提供约 10000 N 以上夹紧力的机构。
拟选用 \textbf{M12 螺栓}作为核心紧固元件。
查阅《机械设计手册》，M12（8.8级高强度螺栓）的标准拧紧力矩 $T \approx 80 N\cdot m$。根据螺纹紧固力的估算公式（简化公式）：
\[ Q \approx \frac{T}{0.2 \times d} \]
代入数据计算螺栓产生的轴向拉力：
\[ Q \approx \frac{80000 \text{ N}\cdot\text{mm}}{0.2 \times 12 \text{ mm}} \approx 33333 \text{ N} \]
\textbf{校验结论}：
即使考虑到使用压板机构时因力臂比（假设为 1:1 或更优）及摩擦导致的力传递损耗，M12 螺栓能提供的实际夹紧力 $W_{act}$ 也远大于理论需求的 $10247$ N（$33333 \text{ N} \gg 10247 \text{ N}$）。
因此，选用 M12 螺旋压板机构不仅满足夹紧要求，还具有较大的安全富余量，结构设计安全可靠。

\subsection{定位误差分析}
本工序需保证孔的位置度。由于使用了 45度斜度夹具，需考虑角度误差。
\begin{enumerate}
    \item 基准不重合误差：$\Delta_B = 0$。
    \item 基准位移误差：$\Delta_Y \approx 0.01$ mm（定位板平面度）。
    \item 夹具角度误差：假设斜度板制造角度误差 $\Delta_\alpha = \pm 1'$，在 100mm 长度上引起的线性误差约为 $100 \times \tan(1') \approx 0.03$ mm。
\end{enumerate}
总误差：
\[ \Delta_D \approx \sqrt{0.01^2 + 0.03^2} \approx 0.032 \text{ mm} \]
满足 IT11 级孔的位置要求（公差通常 $>0.1$ mm）。

\subsection{夹具操作说明}
\begin{enumerate}
    \item 将带斜度的夹具体固定在 X51 铣床工作台上。
    \item 将工件放入斜面定位槽，贴紧底面、侧面及端面挡销。
    \item 旋紧 M12 螺母，压板压紧工件。
    \item 对刀：利用 3mm 塞尺和对刀块，校正主轴 $Y$ 轴及 $Z$ 轴零点。
    \item 启动机床，手动进给钻削 $\Phi14$ 孔，随后更换刀具锪孔。
\end{enumerate}

\section{总结}

回顾设计历程，我经历了一次从理论到实践、从抽象到具象的完整机械加工工艺及装备设计全过程。这对我而言，不仅仅是一次课程作业的完成，更是一次准工程师职业素养的深度洗礼和实战训练。

在设计过程中，我遇到了不少挑战，也收获了宝贵的经验教训。
例如，在工艺规程设计的初期，制定 $\Phi17$ 主轴孔的工艺路线时，我最初并没有深刻意识到 $Ra 0.4 \mu m$ 这一表面粗糙度指标的实现难度，仅草率地安排了“钻-扩-粗铰”的流程。后来，在查阅《机械加工工艺手册》表 3-8 时，我发现粗铰通常只能达到 $Ra 1.6$，这与设计要求相去甚远。这次“险些出错”的经历让我如梦初醒，深刻认识到：工程设计是一项极其严谨的工作，容不得半点想当然和经验主义，任何决策都必须以翔实的数据和国家标准为依据。

此外，在专用夹具的设计环节，我对“定位”与“夹紧”这两个核心概念有了更深层次的理解。在以前的学习中，我往往误以为工件夹紧了就是定位好了。通过这次亲手设计 45 度斜度钻模，我才真正明白：定位的本质是限制自由度，确保工件占据正确的位置；而夹紧的本质是施加外力，保持这个正确位置在加工中不发生变化。两者相辅相成，缺一不可。特别是通过亲手计算切削力，算出看似轻松的钻孔动作竟然伴随着高达 4000 多牛顿的轴向力时，我才意识到夹具强度校验的重要性。如果设计不当，极易导致工件报废甚至引发严重的安全事故。

总的来说，这次课程设计就像一根红线，将我大学期间所学的原本孤立的课程知识串联了起来，形成了一个有机、立体的知识体系。虽然设计中难免还存在稚嫩之处，但我深知自己距离一名优秀的机械工程师还有很长的路要走。这次设计无疑是一个坚实的起点，它培养了我独立思考、解决复杂工程问题的能力，也更加坚定了我投身中国智能制造事业的信心和决心。

\newpage
\section{致谢}

本次课程设计的顺利完成，离不开老师和同学们的每一份帮助与支持。

首先，我要向我的指导教师郭正老师和杨新刚老师致以最崇高的敬意和最诚挚的感谢！在设计选题、方案拟定以及最终成稿的每一个关键环节，老师们都给予了我耐心细致的指导。无论是对工艺难点的点拨，还是对图纸细节的严格把关，老师严谨治学的态度和深厚的学术造诣都让我受益匪浅，令我终身难忘。

其次，我要感谢我的同班同学们。在设计过程中，大家在教室甚至宿舍里共同探讨问题、查阅资料、互相纠错。这种浓厚的学术氛围和团结互助的友谊，是我大学生活中最宝贵的财富。

最后，感谢我的家人，在课程设计期间我身体出现了一些问题，他们给予我无微不至的照顾和支持，让我能够安心完成设计任务。

路漫漫其修远兮，吾将上下而求索。在未来的学习和工作中，我将带着这份感恩之心，继续努力前行！

\begin{thebibliography}{99}
    \bibitem{1} 肖继明, 郑建明. 机械制造技术基础课程设计指导[M]. 北京: 化学工业出版社, 2011.
    \bibitem{2} 王素艳.车床尾座体六面回转钻夹具三维结构设计及有限元分析[J].煤炭技术,2014,33(07):235-237
    \bibitem{3} 李希锋,李琪麟.数控车床自动化夹具定位与夹紧系统设计[J].中国机械,2025,(10):74-76+80.
    \bibitem{4} 尹成湖.机械加工工艺手册[M].化学工业出版社:202301:1014.
\end{thebibliography}

\newpage

\end{document}